% 1-5-structure.tex
%  Budget: 1pp. Describes our own thesis organization -- no citations.
%  Chapter/section names must match the real scaffold (chapter2-5.tex).

\section{ساختار پایان‌نامه}

این پایان‌نامه در پنج فصل سازمان یافته است. فصل نخست \lr{—} همین فصل \lr{—} به
بیان انگیزه و اهمیت موضوع، طرح مسئله، سؤالات پژوهش و نوآوری‌ها و دستاوردهای کار
اختصاص دارد و محدوده و مفروضات کلیِ پژوهش را نیز مشخص می‌کند.

فصل دوم مروری بر پیشینه‌ی پژوهش ارائه می‌دهد؛ کار را با ادبیات پیش‌بینی قیمت برق
از مرورهای نظام‌مند و مطالعات کتاب‌سنجی آغاز می‌کند و سپس به‌ترتیب روش‌های آماری
کلاسیک، یادگیری ماشین کلاسیک، یادگیری عمیق، مدل‌های ترکیبی و مکانیزم توجه،
تفسیرپذیری و هوش مصنوعی توضیح‌پذیر، مطالعات کاربردی و بازارهای خاص، و پیش‌بینی
بلندمدت و احتمالاتی را بررسی می‌کند و با تبیین شکاف پژوهشی و جایگاه کار حاضر به
پایان می‌رسد.

فصل سوم روش تحقیق را تشریح می‌کند: از مفروضات بنیادی و معرفیِ داده‌ها و بازار مورد
مطالعه \lr{—} شامل داده‌ی بنچمارک و داده‌ی زنده \lr{—} و تحلیل اکتشافیِ داده، تا
مهندسی ویژگی، چارچوب اعتبارسنجی و معیارهای ارزیابی، و در نهایت معرفیِ مدل‌ها، از
مدل‌های پایه تا مدل‌های یادگیری ماشین و یادگیری عمیق (\lr{LightGBM} و \lr{LSTM}) و
روش‌های ترکیبی.

فصل چهارم به ارائه‌ی نتایج و تحلیل آن‌ها می‌پردازد؛ شامل نتایج پیش‌بینیِ ساعتی و
روزانه، مقایسه‌ی روش مستقیم و تجمیعی، آزمون معناداریِ آماری (دیبولد-ماریانو) در
برابر مدل‌های این پژوهش و مدل‌های مرجع، تحلیل تفسیرپذیریِ مبتنی بر \lr{SHAP}، و
سرانجام پاسخ صریح به سؤالات پژوهش.

فصل پنجم با جمع‌بندیِ یافته‌ها، بحث درباره‌ی محدودیت‌ها، معرفیِ ابزار
\lr{EnergyMarket-PriceCast} و ارائه‌ی پیشنهادهایی برای پژوهش‌های آینده، به
نتیجه‌گیریِ نهایی می‌رسد.
